% Pursuing Stacks (cotes 134-2 to 134-8): archive cote 134-4, site page 14 (PDF page 15), typescript page 102.
% typescript. Draft compared: ps3.tex:691-761.
% AI-transcribed by Claude Opus 5.5 (claude-opus-5-5, adaptive thinking, effort high), Message Batches API, batch ps1 (msgbatch_01Rv5hCaNTtUt4JHx1yqoHPL), 2026-11-12; image review recorded in records/CHECK_DECISIONS_134_3_4.tsv.
\scanpage{134-4/14}
\mspage{102}
indices.  T\over{H}{h}e dual statement so to say,  when X is decomposed as a sum
$X = \hand{\coprod} X_i$ ,  is valid too,  provided taking products with I is distributive
with respect to \inserted{finite} direct sums.  Thus we get that $M \hand{\to} M_{\underline{I}}$  commutes with \deleted{ar-}%
\deleted{bitrary} \inserted{finite} \xout{\illegible{1}} products,  and with \deleted{arbitrary} \inserted{finite} sums too provided they are dis-%
tributive with respect to multiplication with any object  (or \inserted{with} I \over{only}{only},
which would be enough).

The notion of $\underline{I}$-homotopy in the wide sense between $f,g\colon X \hand{\rightrightarrows} Y$ can
be interpreted in terms of strict $\underline{I}'$-homotopy with variable $\underline{I}'$ ,  as fol-%
lows,  provided we make some mild extra assumptions on (M,W), namely:

a)  (Just for memory)  M is stable under finite products.

b)  M is stable under amalgamated sums  $I \hand{\amalg}_e J$  under the final ob-%
ject e (I and J endowed with sections over e).

c)  If moreover I and J are aspheric over e ,  then so is $I \hand{\amalg}_e J$.

Conditions \over{\unsure{a}}{b}) and \over{\unsure{b}}{c}) give a means of constructing new homotopy inter-%
vals \inserted[typed]{$\underline{K}$},  by amalgamating two \inserted[typed]{homotopy} intervals  $\underline{I}$ and $\underline{J}$ ,  using as sections \xout{for}
of $\underline{I}$ and $\underline{J}$ for making the amalgamation,  either $\hand{\ensuremath{\delta}}_0$ of $\hand{\ensuremath{\delta}}_1$ ,  which gives
four ways of amalgamating  -  of course we take as endpoints of the amal-%
gamat\over{e}{e}\over{d}{d}\xout{\illegible{1}} interval, \inserted[typed]{the} sections over e coming from the two endpoints of
$\underline{I}$  (giving rise to $\hand{\ensuremath{\delta}}_0$ \inserted[typed]{for $\underline{K}$}) and $\underline{J}$  (giving rise to $\hand{\ensuremath{\delta}}_1$ for $\underline{K}$) which have not
been ``used up'' in the amalgamation.  Maybe the handiest convention is to
define \uline{the}  amalgamated interval $\underline{I}\hand{\vee}\underline{J}$ ,  without any ambiguity of choice,
as being
\[ \xout{$\underline{K}$}\ \underline{I}\hand{\vee}\underline{J} = (I,\hand{\ensuremath{\delta}}_1^{\hand{\ensuremath{I}}})\hand{\amalg}_e(J,\hand{\ensuremath{\delta}}_0^{\hand{\ensuremath{J}}}) \qquad \text{endowed with  the two sections coming from } \hand{\ensuremath{\delta}}_0^{\hand{\ensuremath{I}}}\colon e\hand{\to} I \ , \ \hand{\ensuremath{\delta}}_1^{\hand{\ensuremath{J}}}\colon e\hand{\to} J \ , \]
and defining the three other choices \xout{by\illegible{1}} in terms of this operation,  by
replacing one or two among the summands $\underline{I}$, $\underline{J}$  by the ``opposite'' interval
$\overset{\hand{\ensuremath{\vee}}}{\underline{I}}$  or  $\overset{\hand{\ensuremath{\vee}}}{\underline{J}}$ .  T\over[typed]{H}{h}e operation of amalgamation of intervals,  and likewise of
homotopy intervals,  just defined,  is clearly associative up to canonical
isomorphism,  and we have a canonical isomorphism of intervals
\[ (\underline{I}\hand{\vee}\underline{J})^{\hand{\ensuremath{\vee}}} \hand{\simeq} \overset{\hand{\ensuremath{\vee}}}{\underline{I}}\hand{\vee}\overset{\hand{\ensuremath{\vee}}}{\underline{J}} \ . \]
I forgot to check,  for amalgamation of homotopy intervals,  the condition
b) of page 59 namely $e_0 \hand{\cap} e_1 = \hand{\varnothing}_M$  (which has not so far played any role,
anyhow).  To get this condition,  we'll have to be slightly more specific
in condition b) above on M of existence of the relevant amalgamations,
by demanding (as suggested of course by the visual intuition of the situa-%
tion)  that I and \over[typed]{\unsure{U}}{J}  should become subobjects of the amalgamation K,  and
their intersection should be reduce\over{\illegible{1}}{d} to \over[typed]{t}{t}he tautological part e of it.
M$\overset{\hand{\smile}}{\mathrm o}$re relevant still for the use we have in mind is to demand that those
amalgamations should commute to products b\over{y}{y} an arbitrary element X of M.

\margin{14}
